\chapter{MÉTRIQUES DE PERFORMANCE ET MÉTHODES D'ÉVALUATION}
\thispagestyle{empty}

\section{Métriques d'évaluation existantes}

Afin d'assurer la continuité avec la série précédente de rapports, le projet
utilise encore les métriques standard de classement suivantes :
\begin{itemize}
    \item \textbf{Hit@1, Hit@10, Hit@50, Hit@100 :} proportion des cibles de
    test retrouvées parmi les $K$ premières positions classées,
    \item \textbf{MRR :} rang réciproque moyen de la vraie cible sur l'ensemble
    des requêtes d'évaluation,
    \item \textbf{NDCG :} gain cumulatif actualisé normalisé, utilisé pour
    refléter la qualité de la position dans le classement,
    \item \textbf{Recall :} fraction des cibles pertinentes retrouvées dans le
    seuil inspecté,
    \item \textbf{MAP :} précision moyenne moyenne, résumant la qualité du
    classement sur l'ensemble d'évaluation.
\end{itemize}

Ces métriques restent utiles car elles préservent la comparabilité avec les
branches antérieures et quantifient la capacité du système à retrouver des
relations de graphe mises de côté pour le test.

\section{Nouvelles métriques orientées recommandation}

Comme le projet porte en définitive sur la recommandation plutôt que sur la
seule prédiction de liens, la branche révisée ajoute une seconde famille de
métriques :
\begin{itemize}
    \item \textbf{Popularity Lift :} mesure à quel point un modèle se
    concentre sur les articles populaires par rapport à la fréquence de base du
    graphe,
    \item \textbf{Gini Index :} mesure le degré de concentration ou
    d'inégalité de la distribution des recommandations à travers l'ensemble des
    articles,
    \item \textbf{Aggregate Diversity :} compte combien d'articles distincts
    sont effectivement proposés à travers toutes les listes de recommandation,
    \item \textbf{Novelty :} valorise les modèles qui recommandent des articles
    moins évidents,
    \item \textbf{Intra-List Distance (ILD) :} mesure à quel point le contenu
    de chaque liste de recommandation est diversifié.
\end{itemize}

Ces métriques sont importantes car deux modèles ayant des valeurs de Hit@10
similaires peuvent néanmoins se comporter très différemment comme systèmes de
recommandation. Un modèle peut retrouver de nombreuses bonnes relations tout en
proposant sans cesse les mêmes nœuds très connectés ou très populaires. Cette
nouvelle famille de métriques rend ces compromis visibles.

\section{Protocole d'évaluation}

Le protocole officiel d'évaluation actuel est fondé sur la branche structurelle
non orientée révisée. La procédure est résumée dans le
Tableau~\ref{tab:eval-branch-stats}.

\begin{table}[H]
\centering
\caption{Statistiques de la branche d'évaluation révisée et paramètres du protocole.}
\label{tab:eval-branch-stats}
\begin{tabular}{p{0.52\textwidth}p{0.28\textwidth}}
\toprule
\textbf{Quantité} & \textbf{Valeur} \\
\midrule
Arêtes du graphe orienté complet & 291,039 \\
Arêtes structurelles non orientées complètes & 216,123 \\
Arêtes d'entraînement (branche non orientée) & 172,900 \\
Arêtes de test (branche non orientée) & 43,223 \\
Paires d'évaluation étendues & 86,446 \\
Nœuds d'évaluation & 9,022 \\
Coefficient structurel de loi de puissance & $\alpha \approx 2.9966$ \\
Longueur de liste recommandée pour les nouvelles métriques & Top-20 \\
Nombre de modèles dans la comparaison révisée & 10 \\
\bottomrule
\end{tabular}
\end{table}

Dans ce protocole, chaque arête de test non orientée $\{u,v\}$ est étendue en
deux requêtes directionnelles d'évaluation, $u \rightarrow v$ et $v \rightarrow u$,
pour les métriques de classement. Les métriques orientées recommandation sont
ensuite calculées sur les listes Top-20 produites pour chaque nœud
d'évaluation. Cette conception d'évaluation duale permet au projet d'évaluer
à la fois la qualité de récupération du graphe et le comportement pratique de
recommandation au sein d'une même branche expérimentale.
